\documentclass[10pt,letterpaper]{article}
\usepackage[margin=0.85in]{geometry}
\usepackage[T1]{fontenc}
\usepackage{lmodern,microtype,amsmath,amssymb,booktabs,tabularx,array,graphicx}
\usepackage[dvipsnames]{xcolor}
\usepackage{enumitem,fancyhdr}
\usepackage[colorlinks=true,linkcolor=MidnightBlue,urlcolor=MidnightBlue]{hyperref}
\definecolor{accent}{RGB}{20,68,96}
\setlength{\parindent}{0pt}
\setlength{\parskip}{6pt}
\setlength{\tabcolsep}{5pt}
\renewcommand{\arraystretch}{1.18}
\setlist[itemize]{leftmargin=1.4em,itemsep=2pt,topsep=3pt}
\pagestyle{fancy}
\fancyhf{}
\fancyhead[L]{\small\color{accent}SSISDE\quad /\quad Double-well experiment}
\fancyhead[R]{\small Implementation note}
\fancyfoot[C]{\small\thepage}
\setlength{\headheight}{14pt}
\newcommand{\heading}[1]{\vspace{5pt}{\large\bfseries\color{accent}#1}\par}
\newcommand{\file}[1]{\texttt{\detokenize{#1}}}
\newcommand{\RMSE}{\operatorname{RMSE}}
\begin{document}
{\LARGE\bfseries\color{accent}Double well: Black-Scholes workflow}\par
{\small Revised: 10 October 2026}

The OU and double-well notebooks now repeat the Black-Scholes experiment. The extra increment-based drift fits, odd-symmetry fit, and separate OU library have been removed. The earlier long double-well simulation and its results have been replaced. This note supersedes the previous description of those additional comparisons.

\heading{1. Simulation and initial conditions}
Both experiments use $N=100000$, $T=1$, $X_0=1$, and random seed 3, matching Black-Scholes. The fine step is $\Delta t=10^{-5}$. The same Euler simulation convention, time grid, and Brownian indexing are retained.
\begin{align*}
\text{OU:}\qquad dX_t&=-0.5X_t\,dt+0.3\,dB_t^P,\\
\text{Double well:}\qquad dX_t&=(0.5X_t-X_t^3)\,dt+(1+0.25X_t^2)\,dB_t^P.
\end{align*}
The drift differential in the supplied double-well image is interpreted as $dt$. Every drift method receives the known initial drift $\mu(X_0)=-0.5$. ODE2 also receives the known initial spatial slope: $-0.5$ for OU and $-2.5$ for the double well. No slope is estimated from data.

\heading{2. Same experiment stages}
\begin{enumerate}[leftmargin=1.5em,itemsep=5pt]
\item Simulate the process and calculate quadratic variation using the original Black-Scholes scripts' computation and indexing.
\item Estimate $\sigma^2$ with the same degree-4 polynomial library and SR3. Only the diffusion regression parameters are tuned.
\item Calculate diffusion derivatives by the same finite-difference stages as Black-Scholes. The separate analytic-derivative helper is no longer used.
\item Recover $B^Q$, fit the same signature model, and coarse-grain with $M=100$. There are $K=1000$ coarse samples and coarse step $h=0.001$.
\item Compare Method 1, Method 2, and ODE2 with the actual drift. ODE2 uses the corrected sign described below.
\item Recover $B^P$ using \textbf{Method 1 only}, as in Black-Scholes. Method 2 and ODE2 remain drift comparisons.
\item Run \textbf{one} SSISDE fit using that noise path and the original Black-Scholes settings.
\end{enumerate}

The final SSISDE fit has degree-5 drift and diffusion libraries, seven time-series folds, 35 geometrically spaced $\alpha$ values from $10^{-5}$ to $10$, and $\rho\in\{0.3,0.4,0.5,0.7,0.85,0.95,1\}$. The one-standard-error selection rule is unchanged.

There are no additional drift estimators, symmetry assumptions, alternative final libraries, or extra model comparisons. Poor recovery is retained for inspection. All five experiments were rerun after correcting ODE2.

Setting $(\mu/\sigma)''=0$, with derivatives in the state variable, gives
\[
\mu''=2\frac{\sigma'}{\sigma}\mu'
+\left(\frac{\sigma''}{\sigma}-2\frac{(\sigma')^2}{\sigma^2}\right)\mu.
\]
The squared-derivative term previously had a plus sign. Only this sign was corrected; initial conditions and numerical settings are unchanged.

\newpage
\heading{3. Diffusion tuning}
The existing diffusion search fits on the first $80\%$ of the observations, with a three-sample gap before validation on the last $20\%$. Nonconverged candidates and nonpositive predicted variance on observed states are rejected. The selected settings are refitted on the full path. The true diffusion is used only for reporting errors.

OU searches thresholds $\{0.1,0.3,1,3\}$; the double well searches $\{0.01,0.1,0.3,1,3,10\}$. Both search $\nu\in\{0.1,0.5,2\}$, with maximum iterations $10000$ and tolerance $10^{-8}$.

The selected threshold is $3$ in both cases. The selected $\nu$ is $2$ for OU and $0.5$ for the double well. Their fitted variance functions are
\begin{align*}
\hat\sigma^2_{\rm OU}(x)&=0.08940725,\\
\hat\sigma^2_{\rm DW}(x)&=0.98821428+0.02488806x+0.51636376x^2.
\end{align*}
The double-well fit is not forced to contain the exact terms. It is selected from the same polynomial library using observed data.

\heading{4. Completed runs}
The following errors are from the revised notebooks. Diffusion error in the first row refers to the initial SR3 estimate; the final two rows refer to SSISDE.
\begin{center}
\begin{tabular}{@{}lrr@{}}
\toprule
RMSE & OU & Double well\\
\midrule
Initial diffusion & 0.000990 & 0.004308\\
Method 1 drift & 0.192703 & 0.591795\\
Method 2 drift & 0.192586 & 0.423998\\
ODE2 drift & 0.000000 & 1.038716\\
Method 1 Brownian path & 0.320168 & 0.332766\\
Final SSISDE drift & 0.192703 & 0.579673\\
Final SSISDE diffusion & 0.000990 & 0.004688\\
\bottomrule
\end{tabular}
\end{center}
For OU, ODE2's drift RMSE is approximately $2.50\times10^{-8}$ and uses the supplied slope. It does not supply noise to the final SSISDE fit. The final recovered drift is approximately $-0.500311+0.000313962x$ for OU and $-0.523664$ for the double well. These fits do not recover the corresponding true drift functions well, and are left unchanged.

All errors compare against the simulated values using the Black-Scholes workflow's indexing. They are results from one trajectory, not independent test errors. Small SSISDE increment error alone is not proof of accurate drift recovery.

\heading{5. Files and execution}
Run either notebook with \textbf{Run All}:
\begin{itemize}
\item \file{experiments/ou/Stochastic_SINDy_OU.ipynb}
\item \file{experiments/double_well/Stochastic_SINDy_Double_Well.ipynb}
\end{itemize}
The notebooks retain executed plots and printed results. Numerical results are saved in \file{outputs/ou/} and \file{outputs/double_well/}. The ODE2 correction is in \file{models/drift_recovery_from_ODE2.py} and applies to every experiment.
\end{document}
